% چکیده فارسی
\chapter*{چکیده}
\phantomsection
\addcontentsline{toc}{chapter}{چکیده}

\begin{spacing}{1.15}
حل عددی معادلات دیفرانسیل جزئی غیرخطی یکی از مسائل بنیادی در محاسبات علمی و مهندسی به شمار می‌رود. روش‌های عددی کلاسیک با وجود دقت بالا، به تولید شبکه محاسباتی وابسته‌اند و در مواجهه با ابعاد بالا یا هندسه‌های پیچیده با چالش‌های اساسی روبه‌رو می‌شوند. از سوی دیگر، مدل‌های یادگیری عمیق صرفا داده‌محور نیازمند حجم وسیعی از داده‌های تجربی بوده و تضمینی برای ارضای قوانین فیزیکی ارائه نمی‌دهند. چارچوب «شبکه‌های عصبی آگاه از فیزیک» (\lr{PINN}) به‌عنوان رهیافتی نوآورانه برای پیوند دادن معادلات دیفرانسیل حاکم با توانمندی تقریب‌زدن شبکه‌های عصبی عمیق مطرح گردیده است. هدف این پایان‌نامه، مطالعه نظری، پیاده‌سازی محاسباتی و ارزیابی جامع این چارچوب در حل مسائل مستقیم و معکوس معادلات دیفرانسیل جزئی غیرخطی است.

در این پژوهش، پس از تبیین مبانی نظری معادلات دیفرانسیل، روش‌های محاسباتی کلاسیک، یادگیری عمیق و سازوکار مشتق‌گیری خودکار، فرمول‌بندی مدل‌های زمان‌پیوسته و زمان‌گسسته تشریح شده است. در بخش عملی، مدل زمان‌پیوسته برای مسئله معیار معادله غیرخطی برگرز با استفاده از زبان پایتون و کتابخانه محاسباتی جکس پیاده‌سازی گردید. همچنین یک حل‌گر مرجع مستقل مبتنی بر روش طیفی فوریه با گام‌زنی زمانی نمایی مرتبه چهار توسعه یافت و صحت عملکرد آن با جواب نیمه‌تحلیلی تبدیل کول-هوپف راستی‌آزمایی شد.

نتایج حاصل از ارزیابی‌های عددی نشان می‌دهد که در مسئله مستقیم، شبکه عصبی آگاه از فیزیک قادر است تنها با اتکا بر داده‌های محدود مرزی و اولیه و ارضای باقیمانده فیزیکی در نقاط پراکنده، میدان پیوسته پاسخ را در کل دامنه فضا-زمان با دقت بالایی بازسازی نماید. در مسئله معکوس نیز چارچوب پیشنهادی با موفقیت توانست پارامترهای ناشناخته معادله را هم در شرایط داده‌های بدون نویز و هم در حضور اغتشاشات تصادفی اندازه‌گیری با پایداری مناسب شناسایی کند. یافته‌های این پژوهش مؤید کارایی داده‌ای چشمگیر این رهیافت بدون نیاز به تولید شبکه محاسباتی، و انعطاف‌پذیری بالای آن در تحلیل سامانه‌های غیرخطی است.

\vspace{0.3cm}
\noindent\textbf{کلیدواژه‌ها:} شبکه‌های عصبی آگاه از فیزیک، معادلات دیفرانسیل جزئی غیرخطی، مسئله مستقیم، مسئله معکوس، معادله برگرز، مشتق‌گیری خودکار، روش طیفی فوریه.
\end{spacing}
